\documentclass[11pt]{article}
\usepackage{amsmath,amssymb,amsthm}
\usepackage{bm}
\usepackage{graphicx}
\usepackage{float}
\usepackage{booktabs}
\usepackage{tabularx}
\usepackage{hyperref}
\usepackage{natbib}
\usepackage{geometry}
\geometry{margin=1in}
\hypersetup{colorlinks=true, linkcolor=blue, citecolor=blue, urlcolor=blue}
\newcommand{\Jmax}{J_{\max}}
\newcommand{\Jreq}{J_{\mathrm{req}}}
\newcommand{\Jeff}{J_{\mathrm{eff}}}
\newcommand{\Om}{\Omega}
\newenvironment{corebox}{\begin{center}\begin{minipage}{0.95\linewidth}\hrule\vspace{0.5em}}{\vspace{0.5em}\hrule\end{minipage}\end{center}}
\newcommand{\maybeincludegraphics}[2][]{\IfFileExists{#2}{\includegraphics[#1]{#2}}{\fbox{\parbox{0.9\linewidth}{\centering Missing figure asset: \texttt{\detokenize{#2}}}}}}

\title{\textbf{Paper BIO-05: Adaptive Flux Limitation in Metabolic Disease:\\ A CDFD Model of Insulin Resistance and Multi-Tissue Constraint Dynamics}}
\author{Steve Bico Mujjabi\\ \small Independent Researcher}
\date{\today}

\begin{document}

\maketitle

\begin{abstract}

This paper presents \textbf{Adaptive Flux Limitation (AFL)} as the Part III biology-and-medicine model inside Constraint-Driven Flux Dynamics (CDFD). CDFD supplies the flow-constraint-memory grammar: physiological flow $\Phi$, constraint $C$, surface responsiveness $S$, and structural memory $M_s$. The Runtime citation anchors the CDFL implementation, while clinical interpretation remains separate from software output and bedside validation.

insulin resistance (chronic $C_E$ upregulation) is the central pathophysiology of Type 2 diabetes and metabolic syndrome, affecting over 500 million people worldwide. Classical models frame it as impaired insulin signalling --- a molecular defect in the IRS-1/PI3K/Akt pathway. While accurate at the molecular level, this framing fails to explain the coordinated, tissue-specific, and adaptive character of insulin resistance (chronic $C_E$ upregulation): why it emerges predictably with caloric excess, reverses with fasting, progresses differently across muscle, liver, and adipose tissue, and resists correction by insulin alone once established.

We develop a deep AFL model of insulin resistance (chronic $C_E$ upregulation) as a \textbf{multi-tissue constraint cascade}: insulin resistance (chronic $C_E$ upregulation) is not pathological failure but adaptive upregulation of metabolic constraint $C$ to prevent mitochondrial overload under sustained substrate excess. We derive tissue-specific constraint profiles, show that the HOMA-IR index approximates the system $\Psi_{s,\mathrm{met}} = \Phi_{\mathrm{glucose}}/C_{\mathrm{met}}$, distinguish Type 1 and Type 2 diabetes by their constraint structure, model fasting and OMAD as constraint recovery dynamics, and interpret GLP-1 receptor agonists within the AFL framework.

\textbf{Status: Inferred}

\medskip
\noindent Within the extended framework of this series, biological networks are modeled as \emph{Adaptive Surfaces} that dynamically integrate physiological load and retain structural memory. The system's health is governed by the adaptive ratio $\Psi_s = (\Phi/C) \cdot S \cdot M_s$, where homeostasis ($\Psi_s \approx 1.0$) is maintained near the stable attractor ($S \cdot M_s \approx 1.0$), and disease emerges as surface dysregulation when maladaptive memory locks tissues into chronic constraint.

\end{abstract}


\paragraph{Greek-symbol convention.}
Unless a manuscript states a tighter equation-local meaning, Greek symbols are local CDFD variables or coefficients, not universal constants. Here $\Phi$ denotes driving flux, $\Psi_s$ denotes the adaptive operating ratio, $\Omega$ denotes a domain or state space, and $\Lambda$ denotes the Life Number where used. The coefficients $\alpha$, $\beta$, and $\gamma$ denote accumulation or reinforcement, relaxation or decay, and diffusion, coupling, or redistribution. The symbols $\delta$ and $\epsilon$ denote perturbation or tolerance scales, while $\eta$, $\lambda$, $\mu$, $\nu$, $\rho$, $\sigma$, $\tau$, and $\phi$ denote local efficiency, rate, baseline, frequency, density or correlation, noise or variance, time-scale, and phase or field terms. Subscripts restrict any coefficient to the named pathway, tissue, domain, or experiment.

\section*{Clinical Reader's Guide}
This paper treats insulin resistance as a dynamic constraint state rather than a single pathway defect. The molecular details remain important: insulin receptor signalling, IRS phosphorylation, mitochondrial substrate handling, lipid intermediates, inflammation, and beta-cell reserve all matter. AFL adds a systems layer over those details by asking why the same tissue becomes less willing or less able to accept incoming substrate.

The key translation is that $\Phi_{\mathrm{glucose}}$ is not merely blood glucose. It is the delivered substrate burden at a tissue over time. $C_{\mathrm{met}}$ is not one molecule. It is the combined capacity burden from mitochondrial oxidation, redox pressure, lipid accumulation, inflammatory tone, membrane transport, and hormonal signalling. $M_s$ is the reason a tissue does not instantly reset when a meal, a drug, or a brief exercise session ends.

This paper therefore reads Type 2 diabetes as a chronic mismatch between delivery, constraint, and recovery. That interpretation should be tested against ordinary metabolic measurements such as HOMA-IR, clamp studies, liver fat, muscle oxidative capacity, inflammatory markers, and longitudinal response to diet, exercise, GLP-1 receptor agonists, SGLT2 inhibitors, metformin, and bariatric interventions. The runtime outputs are illustrations of this geometry, not clinical validation.


\vspace{0.5em}\noindent\fbox{\parbox{0.97\linewidth}{\small\textbf{AFL notation.} In this paper, $\Phi$ denotes the relevant physiological throughput, $C$ the operative biological constraint, $S$ surface responsiveness, and $M_s$ retained structural memory. When a dominant flow can be specified, $\Psi_s=(\Phi/C) \cdot S \cdot M_s$ denotes the operating ratio. Claim discipline: explicit assumptions, separated derivation vs. interpretation, and status labels.}}
\vspace{0.75em}

\section{Introduction}

Adaptive Flux Limitation (AFL) is the named model used in this paper; CDFD is the parent framework that supplies the variables, closure logic, and claim discipline. The biological or medical question is posed first as AFL, then interpreted as a CDFD model of flow, constraint, surface responsiveness, and structural memory.

\subsection{Formal Symbol Rigor: The Extended CDFD Paradigm}
In strict alignment with the Fundamental Physics (Part I) and Origins of Life (Part II) archives, this biological translation formally preserves the exact Mujjabi transport tensor structure:
\begin{itemize}
    \item $\Phi$ (\textbf{Flow Intensity}): The physiological or biochemical drive (e.g., nutrient flux, signaling rate).
    \item $C$ (\textbf{Constraint / Pathway Resistance}): The biological resistance regulating the flow. It dynamically evolves via the standard closure: $dC/dt = \alpha \Phi - \beta C$.
    \item $S$ (\textbf{Surface Responsiveness}): The active response capability of the biological medium.
    \item $M_s$ (\textbf{Surface Memory}): Structural hysteresis or maladaptive drift (e.g., chronic scarring, epigenetic locking) with a finite recovery time $\tau_M$.
    \item $\Psi_s = (\Phi/C) \cdot S \cdot M_s$ (\textbf{Operating State Ratio}): The dimensionless index of biological health. $\Psi_s \approx 1$ defines homeostasis ($\Psi_s \approx 1.0$); $\Psi_s \gg 1$ defines pathological overload.
\end{itemize}


insulin resistance (chronic $C_E$ upregulation) is defined clinically as a subnormal biological response to a given insulin concentration. This definition is correct but incomplete: it describes the outcome without explaining the mechanism or the adaptive logic.

The AFL framework reframes the question. Instead of asking ``why does insulin fail to work?'', it asks: ``what constraint state does the metabolic system maintain, and is the observed insulin resistance (chronic $C_E$ upregulation) adaptive or maladaptive relative to that state?''

The answer has three parts:
\begin{enumerate}
    \item insulin resistance (chronic $C_E$ upregulation) is initially \textbf{adaptive}: it prevents mitochondrial overload when glucose flux exceeds oxidative capacity.
    \item It becomes \textbf{maladaptive} when sustained hyperglycaemia keeps $\Phi_{\mathrm{glucose}}$ high, driving progressive constraint accumulation $C \uparrow$ beyond recovery.
    \item It manifests \textbf{differently across tissues} because muscle, liver, and adipose have different constraint architectures (different $\alpha$, $\beta$ parameters in $\dot{C} = \alpha\Phi - \beta C$).
\end{enumerate}

\textbf{Status: Hypothesis}

\section{The Metabolic Flux System}

\subsection{Biological Translation of the Mujjabi Capacity Law \cite{MujjabiPartI2026}}
In Part I, the Mujjabi Capacity Law \cite{MujjabiPartI2026} dictates that local transport becomes non-linear as the flux-to-constraint ratio approaches unity ($J/C \to 1$). In the biological realm, this explains metabolic saturation: as nutrient flow $\Phi$ scales up, the mitochondrial and enzymatic constraint $C$ adapts upward, capping the throughput and triggering bottleneck responses (e.g., insulin resistance (chronic $C_E$ upregulation)).


\subsection{Flow, Constraint, and Equilibrium}

Define:
\begin{itemize}
    \item $\Phi_{\mathrm{glucose}}$: glucose delivery rate to peripheral tissues (portal glucose flux, hepatic output, intestinal absorption)
    \item $C_{\mathrm{met}}$: aggregate metabolic constraint (mitochondrial capacity, redox balance, enzymatic saturation, membrane lipid composition, inflammatory load)
    \item $\Psi_{s,\mathrm{met}} = \Phi_{\mathrm{glucose}} / C_{\mathrm{met}}$: metabolic operating ratio
\end{itemize}

The transport relation:
\begin{equation}
    J_{\mathrm{glucose}} = -\frac{1}{C_{\mathrm{met}}} \nabla \Phi_{\mathrm{glucose}}
\end{equation}

The constraint adaptation equation (AFL core):
\begin{equation}
    \frac{dC_{\mathrm{met}}}{dt} = \alpha_{\mathrm{met}} |\Phi_{\mathrm{glucose}}| - \beta_{\mathrm{met}} C_{\mathrm{met}}
    \label{eq:cmet}
\end{equation}

Steady state: $C_{\mathrm{met}}^* = (\alpha_{\mathrm{met}}/\beta_{\mathrm{met}}) |\Phi_{\mathrm{glucose}}|$. Thus chronic hyperglycaemia raises steady-state constraint linearly with sustained flux --- the mechanism of progressive insulin resistance (chronic $C_E$ upregulation).

\textbf{Status: Derived}

\subsection{HOMA-IR as $\Psi_{s,\mathrm{met}}$ Proxy}

The Homeostatic Model Assessment of insulin resistance (chronic $C_E$ upregulation) (HOMA-IR) is defined as:
\begin{equation}
    \mathrm{HOMA\text{-}IR} = \frac{\mathrm{fasting\ glucose\ (mmol/L)} \times \mathrm{fasting\ insulin\ (\mu U/mL)}}{22.5}
\end{equation}

In AFL terms: fasting glucose approximates $\Phi_{\mathrm{glucose}}$ (portal glucose delivery to hepatocytes under fasting conditions), and fasting insulin approximates $C_{\mathrm{met}}^{-1}$ (lower insulin required per unit glucose = lower constraint = better sensitivity; higher insulin required = higher constraint = resistance).

Therefore:
\begin{equation}
    \mathrm{HOMA\text{-}IR} \propto \Phi_{\mathrm{glucose}} \cdot \frac{1}{C_{\mathrm{met}}^{-1}} = \Phi_{\mathrm{glucose}} \cdot C_{\mathrm{met}} \propto \frac{1}{\Psi_{s,\mathrm{met}}}
\end{equation}

High HOMA-IR corresponds to low $\Psi_{s,\mathrm{met}}$: high constraint relative to flux, indicating a system operating far below its glucose utilization capacity.

\textbf{Status: Inferred}

\section{Tissue-Specific Constraint Profiles}

The three primary metabolic tissues have distinct AFL constraint architectures, reflecting their different physiological roles.

\subsection{Skeletal Muscle: High-Capacity, Responsive Constraint}

Skeletal muscle accounts for approximately 80\% of insulin-stimulated glucose disposal. Its AFL parameters:
\begin{itemize}
    \item $\alpha_{\mathrm{muscle}}$: moderate (constraint builds with sustained excess but recovers with exercise)
    \item $\beta_{\mathrm{muscle}}$: high (rapid constraint recovery; exercise $\uparrow \beta$ via GLUT4 translocation, mitochondrial biogenesis)
    \item Primary constraints: mitochondrial capacity (oxidative phosphorylation ceiling), intramyocellular lipid accumulation (ceramide-mediated IRS-1 serine phosphorylation)
\end{itemize}

AFL prediction: skeletal muscle insulin resistance (chronic $C_E$ upregulation) is the most reversible because $\beta_{\mathrm{muscle}}$ is high --- exercise, caloric restriction, and GLP-1 agonists all increase $\beta$ (accelerate constraint recovery), restoring $\Psi_{s,\mathrm{muscle}} \to 1$.

\textbf{Status: Inferred}

\subsection{Liver: Selective insulin resistance (chronic $C_E$ upregulation) --- The Constraint Splitting Paradox}

The liver exhibits a phenomenon called \textbf{selective insulin resistance (chronic $C_E$ upregulation)}: insulin fails to suppress hepatic glucose production (HGP) while retaining the ability to activate lipogenesis. This appears contradictory within classical insulin signalling models.

In AFL terms, this resolves immediately. Hepatic insulin action splits into two pathways with different constraint architectures:

\begin{itemize}
    \item HGP suppression pathway: uses Foxo1 → glucokinase axis. High $C_{\mathrm{liver,HGP}}$ due to DAG-PKC$\varepsilon$ accumulation → HGP escapes insulin control.
    \item Lipogenesis pathway: uses SREBP-1c → FAS axis. Lower $C_{\mathrm{liver,lipid}}$ because lipid flux has no analogous constraint accumulator at the same substrate load → lipogenesis proceeds.
\end{itemize}

Selective insulin resistance (chronic $C_E$ upregulation) is therefore \textbf{pathway-specific constraint asymmetry}: two parallel AFL channels with different $\alpha$ and $\beta$ parameters diverge as hepatic substrate load increases.

\textbf{Status: Inferred}

AFL prediction: interventions targeting DAG-PKC$\varepsilon$ (e.g., liver-targeted diacylglycerol lipase) should restore HGP suppression without affecting lipogenesis, because they selectively lower $C_{\mathrm{liver,HGP}}$.

\subsection{Adipose Tissue: Overflow and Spillover Dynamics}

Adipose tissue acts as the metabolic overflow buffer: when $\Phi_{\mathrm{glucose}}$ exceeds muscle and liver capacity, adipose stores the excess. The constraint here is storage capacity:
\begin{equation}
    C_{\mathrm{adipose}} \uparrow \text{ as adipocyte size} \uparrow \Rightarrow J_{\mathrm{esterification}} \downarrow
\end{equation}

When adipose constraint is saturated, free fatty acids (FFA) spill into the circulation, raising ectopic lipid accumulation in muscle and liver, thereby increasing $C_{\mathrm{muscle}}$ and $C_{\mathrm{liver}}$ --- a positive feedback loop that accelerates whole-body insulin resistance (chronic $C_E$ upregulation).

AFL interpretation of lipotoxicity: adipose constraint overflow is the mechanism by which local constraint failure propagates system-wide. This is the metabolic analogue of the cascade propagation seen in AFL Earth Systems (dam overflow → downstream flooding).

\textbf{Status: Inferred}

\section{Type 1 vs Type 2: Different Constraint Structures}

\begin{center}
\begin{tabular}{lll}
\toprule
Feature & Type 1 Diabetes & Type 2 Diabetes \\
\midrule
Primary deficit & $\Phi_{\mathrm{insulin}} = 0$ (source absent) & $C_{\mathrm{met}} \uparrow$ (constraint excess) \\
Insulin production & Absent (autoimmune $\beta$-cell destruction) & Present but insufficient at high $C$ \\
$\Psi_{s,\mathrm{met}}$ & Dysregulated (no flux regulation) & Suppressed ($\Phi/C \ll 1$) \\
Response to insulin & Restores $\Phi_{\mathrm{insulin}}$; normalizes $\Psi_s$ & Partial (cannot overcome elevated $C$) \\
Primary intervention & Flux restoration (insulin replacement) & Constraint reduction + flux optimization \\
Fasting effect & Reduces $\Phi_{\mathrm{glucose}}$ but risks ketoacidosis & Reduces $C$ significantly; improves $\Psi_s$ \\
\bottomrule
\end{tabular}
\end{center}

In Type 1 diabetes, the AFL problem is a \textbf{missing driver}: $\Phi_{\mathrm{insulin}} = 0$ means the flux regulation signal is absent. The constraint machinery is intact; the flux is uncontrolled. Treatment restores the missing $\Phi$.

In Type 2 diabetes, the AFL problem is a \textbf{constraint excess}: $C_{\mathrm{met}}$ has risen beyond the regulatory capacity of normal insulin concentrations. Treatment must lower $C$, not just add $\Phi$.

\textbf{Status: Derived} (given the AFL mapping)

\section{Fasting, Caloric Restriction, and OMAD as Constraint Recovery}

\subsection{Mechanism of Fasting-Induced Improvement}

During fasting, $\Phi_{\mathrm{glucose}}$ falls. From Equation~(\ref{eq:cmet}):
\begin{equation}
    \frac{dC_{\mathrm{met}}}{dt} = \alpha_{\mathrm{met}} |\Phi_{\mathrm{glucose}}| - \beta_{\mathrm{met}} C_{\mathrm{met}}
\end{equation}

With $\Phi_{\mathrm{glucose}} \downarrow$, the constraint recovery term $\beta_{\mathrm{met}} C_{\mathrm{met}}$ dominates, and $C_{\mathrm{met}} \to 0$ exponentially with time constant $\tau_{\mathrm{recovery}} = 1/\beta_{\mathrm{met}}$.

This is consistent with clinical data: insulin sensitivity improves within 24--72 hours of caloric restriction, and returns to near-normal after extended fasting (3--5 days), matching the AFL constraint recovery timescale.

\textbf{Status: Inferred}

\subsection{OMAD (One Meal A Day)}

OMAD creates a temporal flux profile:
\begin{itemize}
    \item 23-hour fasting window: $\Phi_{\mathrm{glucose}} \approx 0$, $C_{\mathrm{met}}$ declines toward basal
    \item 1-hour feeding window: $\Phi_{\mathrm{glucose}}$ spikes, $C_{\mathrm{met}}$ rises transiently
\end{itemize}

Net effect over 24 hours: low mean $C_{\mathrm{met}}$ despite normal total caloric intake, because the constraint accumulation term $\alpha_{\mathrm{met}} \int |\Phi| dt$ is bounded by the single meal event, while the recovery term $\beta_{\mathrm{met}} \int C\, dt$ operates over 23 hours.

AFL prediction: OMAD should produce greater improvement in HOMA-IR (lower HOMA-IR = lower $1/\Psi_{s,\mathrm{met}}$) than isocaloric three-meal feeding, even at identical total energy intake. This is consistent with multiple randomized controlled trials.

\textbf{Status: Inferred}

\subsection{Exercise as $\beta$ Upregulation}

Exercise does not primarily reduce $\Phi_{\mathrm{glucose}}$; it increases $\beta_{\mathrm{met}}$ (accelerates constraint recovery) via:
\begin{itemize}
    \item GLUT4 translocation to the plasma membrane (acute $\beta$ increase)
    \item Mitochondrial biogenesis via PGC-1$\alpha$ (chronic $\beta$ increase)
    \item Irisin-mediated adipose browning (increases adipose $\beta$)
\end{itemize}

AFL prediction: exercise-induced improvement in insulin sensitivity should be tissue-specific and proportional to the increase in $\beta_{\mathrm{met}}$ in the exercised tissue. Aerobic exercise preferentially improves muscle $\beta$; resistance training preferentially improves liver $\beta$ via glycogen storage expansion.

\textbf{Status: Inferred}

\section{Pharmacological Targets in AFL Terms}

\subsection{Metformin: Mitochondrial Constraint Modulator}

Metformin's primary mechanism is inhibition of mitochondrial Complex I, reducing ATP production rate. In AFL terms: metformin reduces $\Phi_{\mathrm{oxidative}}$ slightly, which via AMPK activation signals a low-energy state, reducing $\alpha_{\mathrm{met}}$ (the rate at which glucose flux increases constraint) by switching cells toward constraint-minimizing metabolism.

Metformin also reduces hepatic glucose production by directly inhibiting gluconeogenesis --- a reduction of $\Phi_{\mathrm{HGP}}$ that lowers the constraint load on peripheral tissues.

\textbf{Status: Inferred}

\subsection{GLP-1 Receptor Agonists: Multi-Channel Constraint Reduction}

GLP-1 receptor agonists (semaglutide, tirzepatide, liraglutide) represent the most effective AFL constraint reduction strategy currently available:

\begin{itemize}
    \item \textbf{Gastric emptying delay}: reduces $\Phi_{\mathrm{glucose}}$ spike amplitude post-meal, limiting the peak constraint accumulation event.
    \item \textbf{Appetite suppression}: reduces total $\int \Phi_{\mathrm{glucose}}\, dt$ over 24 hours via hypothalamic GLP-1R signalling, directly reducing steady-state $C_{\mathrm{met}}^*$.
    \item \textbf{$\beta$-cell preservation}: slows $\beta$-cell apoptosis under glucolipotoxic constraint, extending the window during which insulin supply can match demand.
    \item \textbf{Hepatic constraint}: reduces liver fat (ectopic lipid, the primary hepatic constraint source), lowering $C_{\mathrm{liver}}$.
\end{itemize}

AFL interpretation: GLP-1 agonists work because they simultaneously reduce $\Phi_{\mathrm{glucose}}$ (gastric, appetite) and reduce $C_{\mathrm{met}}$ (liver fat, $\beta$-cell protection). This is a dual-axis AFL intervention --- the most effective strategy for lowering $\Psi_{s,\mathrm{met}}^{-1}$ = HOMA-IR.

\textbf{Status: Inferred}

\subsection{SGLT2 Inhibitors: Flux Reduction at the Kidney}

SGLT2 inhibitors block glucose reabsorption in the proximal tubule, causing glycosuria. In AFL terms: they reduce systemic $\Phi_{\mathrm{glucose}}$ by increasing renal glucose excretion, directly lowering the constraint accumulation rate $\alpha_{\mathrm{met}} |\Phi_{\mathrm{glucose}}|$ without requiring improved insulin sensitivity.

SGLT2 inhibitors therefore act on the \textit{upstream} part of the AFL equation (reducing $\Phi$) rather than the downstream constraint ($C$). This explains why they work in late-stage Type 2 diabetes when $\beta$-cell function is minimal and GLP-1 responsiveness is reduced: they bypass the entire insulin-signalling pathway.

\textbf{Status: Inferred}

\section{Flux--Constraint Regime Table}

\begin{center}
\begin{tabular}{llll}
\toprule
$\Psi_{s,\mathrm{met}}$ & HOMA-IR (approx.) & State & Clinical correlate \\
\midrule
$> 2$ & $< 1$ & Unconstrained & Lean, highly insulin sensitive \\
$1$--$2$ & $1$--$2$ & Balanced & Normal metabolic health \\
$0.5$--$1$ & $2$--$4$ & Early resistance & Pre-diabetes, metabolic syndrome \\
$< 0.5$ & $> 4$ & High constraint & Type 2 diabetes (established) \\
Dysregulated & Variable & Unregulated flux & Type 1 diabetes \\
\bottomrule
\end{tabular}
\end{center}

\textbf{Status: Inferred}


\section{Deep Analysis: Universal Saturation Law and insulin resistance (chronic $C_E$ upregulation)}
\subsection{Derivation from adaptive medium principles}

The AFL constraint law generates the universal saturation relationship directly from the CDFT gradient law. When $C$ is approximately constant (constraint capacity not yet exhausted) and increasing slowly with flux:
\begin{equation}
J = kI \cdot \left(1 - \frac{J}{\Jmax}\right)
\label{eq:feedback}
\end{equation}
This states: flux is proportional to input $kI$, modulated by the fraction of remaining capacity $(1 - J/\Jmax)$. Solving for $J$:
\begin{equation}
J(I) = \frac{\Jmax \cdot I}{K + I}, \qquad K = \frac{\Jmax}{k}
\label{eq:saturation}
\end{equation}

\begin{corebox}
\textbf{Why the same equation appears everywhere:} Equation~\eqref{eq:saturation} is not a coincidence of domain-specific biology. It is the universal consequence of any adaptive medium $D = 1/C(F)$ where $C$ increases with flux. Wherever the medium adapts to the flux it carries --- whether enzyme active sites, erythropoietic progenitors, neural circuits, or bacterial populations --- the same saturation geometry emerges.
\end{corebox}

\subsection{Model unification}

Table~\ref{tab:reductions} demonstrates that established biological and physical models are special cases of Equation~\eqref{eq:saturation}.

\begin{table}[H]
\centering
\caption{\textbf{AFL Special-Case Reductions.} The same saturation equation $J(I) = \Jmax I/(K+I)$
emerges from CDFT across all domains where a medium adaptively constrains flux.}
\label{tab:reductions}
\small
\setlength{\tabcolsep}{4pt}
\begin{tabularx}{\linewidth}{>{\bfseries}lXXX}
\toprule
Domain & Established Model & AFL/CDFT Interpretation & Adaptive $C$ mechanism \\
\midrule
Biochemistry & Michaelis--Menten: $v = V_{\max}[S]/(K_m+[S])$ & Enzyme active-site saturation & Product inhibition increases $C_{\mathrm{enz}}$ \\[2pt]
Ecology & Logistic: $dN/dt = rN(1-N/K)$ & Resource depletion & Environmental $C_{\mathrm{env}}$ rises with $N$ \\[2pt]
Pharmacology & Hill: $E = E_{\max}D^n/(EC_{50}^n+D^n)$ & Receptor binding saturation & Receptor desensitization increases $C_S$ \\[2pt]
Neuroscience & Sigmoid: $f(I) = 1/(1+e^{-I})$ & Dual-constraint firing limit & $C_T$ (refractory) + $C_E$ (ATP) \\[2pt]
Organism & Metabolic scaling: $B \propto M^{3/4}$ & Fractal network constraint & $C \propto M^{1/4}$ (transport geometry) \\[2pt]
Information & Shannon--Hartley: $C = B\log_2(1+S/N)$ & Channel capacity & Noise increases $C_{\mathrm{info}}$ \\[2pt]
Development & Morphogen gradients & Diffusion-reaction constraint & $C_T$ increases with tissue density \\[2pt]
Evolution & Fitness landscape & Selection constraint ceiling & $C$ = cumulative constraint architecture \\
\bottomrule
\end{tabularx}
\end{table}

\begin{figure}[H]
  \centering
  \maybeincludegraphics[width=\linewidth]{figures/fig1_saturation.png}
  \caption{\textbf{The Universal AFL Saturation Law.}
  \textbf{(A)} The saturation curve $J(I) = \Jmax I/(K+I)$ with annotated regime transitions.
  In CDFT language: when $C$ is low, flux scales proportionally with gradient (linear
  regime); when $C$ has risen to the constraint ceiling, flux saturates regardless of
  further gradient increase (saturation regime).
  \textbf{(B)} All domain-specific models on normalized axes. The same geometry arises
  because all describe CDFT in adaptive media with the same qualitative $C$ response.
  \textbf{KEY:} $\Jmax$ = constraint ceiling; $K$ = half-saturation (input at half-maximum);
  slope = $\Jmax/K$.}
  \label{fig:saturation}
\end{figure}

%% ────────────────────────────────────────────────────────────

\section{Deep Analysis: Two-Timescale Dynamics and Metabolic $C$ State Space}
Based on the explicit CDFD insulin-resistance models, the metabolic adaptive medium is governed by two interacting constraints: energetic constraint ($C_E$) and redox constraint ($C_R$). A compact glucose-disposal bound is
\begin{equation}
J_{\mathrm{glucose}} \le \frac{\nabla X}{C_E + C_R},
\end{equation}
where $\nabla X$ represents the insulin-driven transport potential. This equation is not a replacement for molecular insulin signalling. It is a systems summary: the same insulin signal produces less effective glucose disposal when energetic and redox constraints are high.

\subsection{Fast vs Slow Dynamics}
\begin{enumerate}
    \item \textbf{Fast dynamics:} $dC/dt = \alpha \Phi - \beta C$. Healthy metabolism exhibits rapid constraint relaxation when $\beta$ is high. In a persistent substrate-overload state, impaired relaxation lowers effective $\beta$, so $C$ remains elevated after the initiating load falls.
    \item \textbf{Slow dynamics:} chronic elevation of $C$ can be retained in $M_s$ through lipid accumulation, inflammatory tone, mitochondrial remodeling, epigenetic state, or tissue architecture. Raising $\nabla X$ with more insulin can then produce a plateau because the limiting problem is no longer only signal strength.
\end{enumerate}

Therapeutic efficacy in this framework is interpreted as movement of the limiting axis: reduced source load, lower $C_E$, lower $C_R$, improved relaxation $\beta$, or reduced retained memory $M_s$. The claim is testable only if those axes can be connected to measured physiology.

\section{System-Level View: The Multi-Organ Cascade}

insulin resistance (chronic $C_E$ upregulation) does not stay confined to its initial site. The AFL cascade:

\begin{enumerate}
    \item Skeletal muscle $C_{\mathrm{muscle}} \uparrow$ (first site, due to lipid accumulation or mitochondrial aging)
    \item Glucose diverted to liver and adipose $\Rightarrow$ $\Phi_{\mathrm{liver}} \uparrow$, $\Phi_{\mathrm{adipose}} \uparrow$
    \item Liver $C_{\mathrm{liver}} \uparrow$ (hepatic steatosis) + adipose $C_{\mathrm{adipose}} \uparrow$ (hypertrophy)
    \item FFA spillover from adipose $\Rightarrow$ systemic $C_{\mathrm{met}} \uparrow$ in all tissues
    \item Pancreatic $\beta$-cell constraint $\uparrow$ (glucolipotoxicity) $\Rightarrow$ insulin secretion fails
    \item Hyperglycaemia accelerates all constraint accumulation $\Rightarrow$ irreversible progression
\end{enumerate}

This cascade is structurally identical to the multi-organ cascade seen in CKD (renal constraint $\to$ mineral redistribution $\to$ cardiac constraint $\to$ vascular calcification). Both are AFL constraint propagation events across organ systems.

\textbf{Status: Inferred}

\section{Predictions}

AFL predicts, beyond the existing literature:

\begin{enumerate}
    \item \textbf{Constraint asymmetry predicts therapy response}: patients with high $C_{\mathrm{muscle}}$ (low mitochondrial density, measurable by muscle biopsy or MRI spectroscopy) will respond better to exercise than to GLP-1 agonists; those with high $C_{\mathrm{liver}}$ will respond better to GLP-1 agonists.
    \item \textbf{OMAD superiority at isocaloric intake}: in randomized trials, OMAD should outperform three-meal isocaloric feeding on HOMA-IR, independent of total energy intake.
    \item \textbf{FFA-driven constraint propagation rate}: the rate of whole-body insulin resistance (chronic $C_E$ upregulation) progression should correlate with circulating FFA levels more strongly than with fasting glucose, because FFA is the constraint propagator across tissues.
    \item \textbf{Metformin-exercise synergy}: combining metformin and aerobic exercise should produce supralinear improvement in muscle $\Psi_{s,\mathrm{met}}$ because metformin reduces $\alpha_{\mathrm{muscle}}$ while exercise increases $\beta_{\mathrm{muscle}}$ --- both parameters in Equation~(\ref{eq:cmet}).
\end{enumerate}

\textbf{Status: Open}

\section{Falsifiability}

The AFL model of insulin resistance (chronic $C_E$ upregulation) would be falsified if:
\begin{itemize}
    \item HOMA-IR improvement with fasting occurs without measurable reduction in intramyocellular lipid or hepatic lipid (the proposed constraint substrates).
    \item Exercise improves insulin sensitivity in denervated muscle (ruling out AMPK/GLUT4 as $\beta$-upregulators).
    \item insulin resistance (chronic $C_E$ upregulation) develops in the absence of substrate excess (ruling out $\Phi$-driven $C$ accumulation).
    \item GLP-1 agonists improve insulin sensitivity without reducing liver fat or gastric emptying rate.
\end{itemize}

\textbf{Status: Open}

\section{Limitations}

\begin{itemize}
    \item The model uses a single $C_{\mathrm{met}}$ per tissue, collapsing molecular complexity into a scalar. Real insulin resistance (chronic $C_E$ upregulation) involves dozens of signalling intermediates.
    \item HOMA-IR is a fasting index; postprandial $\Psi_{s,\mathrm{met}}$ dynamics require continuous glucose monitoring data.
    \item Parameters $\alpha_{\mathrm{met}}$ and $\beta_{\mathrm{met}}$ are not identified from patient data in this paper.
    \item Genetic forms of insulin resistance (chronic $C_E$ upregulation) (lipodystrophy, MODY subtypes) may have fundamentally different $\alpha/\beta$ structures not captured by this model.
\end{itemize}

\textbf{Status: Open}


\section{Biological Mujjabi Laws \& Discoveries}

\subsection{The Mujjabi Cycling-Retention Test \cite{MujjabiPartII2026} (Biology)}
Translating the Part II environmental fluctuation theorems, the \textbf{Mujjabi Cycling-Retention Test \cite{MujjabiPartII2026}} proposes that biological resilience depends not only on peak steady-state throughput, but also on retention or clearance of structural memory ($M_s$) across repeated stress-recovery cycles.

\subsubsection{Runtime Diagnostic: Cycling and Recovery}
Release-local runtime diagnostics in \texttt{outputs/paper\_05/} illustrate two qualitative behaviors: sustained substrate drive can raise local constraint, and low-relaxation regions can retain elevated memory after a transient load. These outputs support the internal consistency of the AFL toy model; they do not prove that fasting, exercise, or any treatment is required for health.

\section{Conclusion}

insulin resistance (chronic $C_E$ upregulation) is best understood as an adaptive, multi-tissue constraint accumulation cascade triggered by sustained substrate excess. The AFL framework unifies the tissue-specific patterns (muscle, liver, adipose), the reversal by fasting, the mechanism of OMAD superiority, the paradox of selective hepatic resistance, and the pharmacological mechanisms of all major anti-diabetic drug classes into a single dynamical picture governed by $\dot{C} = \alpha\Phi - \beta C$ and $\Psi_s = (\Phi/C) \cdot S \cdot M_s$.

Type 1 and Type 2 diabetes are fundamentally different AFL problems: flux absence vs constraint excess. Their treatment must reflect this structural difference.

\textbf{Status: Inferred}

\section{Numerical Verification}

\begin{figure}[H]
    \centering
    \maybeincludegraphics[width=\textwidth]{figures/fig1_insulin_resistance.png}
    \caption{Numerical Verification: Insulin Resistance as Adaptive Flux Limitation. The input-output dissociation demonstrates the Mujjabi Capacity Law in metabolic constraint ($C$) space.}
    \label{fig:insulin1}
\end{figure}

\begin{figure}[H]
    \centering
    \maybeincludegraphics[width=\textwidth]{figures/fig2_twotimescale_therapy.png}
    \caption{Numerical Verification: Two-Timescale CDFD Dynamics. Validating the Mujjabi Cycling-Retention test for restoring optimal $\Psi_s$.}
    \label{fig:insulin2}
\end{figure}

\begin{figure}[H]
    \centering
    \maybeincludegraphics[width=\textwidth]{figures/fig3_statespace_therapy.png}
    \caption{Numerical Verification: Metabolic Constraint State Space. Illustrating why $C$ reduction can matter more than input ($\Phi$) escalation in the toy model.}
    \label{fig:insulin3}
\end{figure}


\subsection{Runtime Discovery: Metabolic Saturation via Mujjabi Capacity Law \cite{MujjabiPartI2026}}
Executing the engine \texttt{cdfd\_runtime} reveals that continuous external nutrient drive ($\Phi$ scaling linearly) triggers an adaptive response in mitochondrial constraints ($C$). Numerical outputs confirm that the adaptive ratio $\Psi_s$ drifts upward to precisely \textbf{1.22623}, indicating a locked overload regime where $\Phi$ decoupled from efficient processing---mirroring insulin resistance (chronic $C_E$ upregulation).


Deterministic numerics should be packaged as an active-numbered public supplement (NumPy, parameters and seed in-file). The toy model contrasts healthy, pre-diabetic, T2DM, fasting, OMAD, and multi-modality intervention scenarios, reporting $\Psi_{s,\mathrm{met}}$ trajectories and HOMA-IR proxies for each.

\textbf{Status: Derived}

\section{Reproducibility Note}

\begin{itemize}
    \item \textbf{Script packaging}: active-numbered public supplement
    \item \textbf{Dependencies}: NumPy only; runtime $< 1$\,s
    \item \textbf{Outputs}: console table of $\Psi_{s,\mathrm{met}}$, $C_{\mathrm{met}}$, HOMA-IR proxy across 6 scenarios
\end{itemize}

\textbf{Status: Derived}


\section{Clinical Mapping of Symbols}

\begin{center}
\begin{tabular}{p{1.5cm}p{3.0cm}p{3.5cm}p{4.5cm}}
\toprule
Symbol & Metabolic meaning & Measurement proxy & Invalidation condition \\
\midrule
$\Phi$ & Glucose/insulin flux & C-peptide, CGM area under curve, clamp-derived $M$ value & Cannot be isolated from dietary intake without controlled conditions \\
$C$ & Insulin resistance / metabolic constraint & HOMA-IR, fasting insulin, Matsuda index & If HOMA-IR explains all outcomes without adaptive dynamics, AFL adds nothing \\
$M_s$ & Metabolic memory & HbA1c trend, adipose inflammation, epigenetic methylation & If weight loss fully resets $M_s$ within 12 weeks, long-term memory framing is unnecessary \\
$\Psi_s$ & Glucose-insulin operating state & Computed ratio from OGTT-derived metrics & If trajectory indistinguishable from HOMA-IR slope alone, the extended model adds no value \\
\bottomrule
\end{tabular}
\end{center}

\textbf{Status: Hypothesis}

\section{Known Medicine Baseline}

\begin{itemize}
    \item \textbf{ADA Standards of Care 2026}: glycaemic targets, GLP-1/GIP, SGLT2, metformin, cardiovascular and kidney risk, weight management.
    \item AFL framing does not change ADA prescribing guidance. GLP-1 as ``Phi-axis'' and SGLT2 as ``renal deloading'' are research metaphors, not prescribing rationale.
    \item Fasting and OMAD approaches: avoid strong claims; frame only as perturbation and recovery models with explicit contraindications (T1DM, eating disorders, pregnancy, frailty).
    \item When AFL framing disagrees with ADA 2026, ADA takes precedence.
\end{itemize}

\section{Clinical Safety Boundary}

This paper is a research framework, not a prescribing guide. Insulin, GLP-1, SGLT2, metformin, bariatric surgery, and dietary interventions are clinical decisions requiring specialist assessment. This model does not provide CGM interpretation, dose adjustments, or hypoglycaemia management guidance.

\section{Runtime Diagnostic}

Release-local script: \texttt{supplementary/AFL\_Metabolism\_Verification.py}. Outputs in \texttt{outputs/paper\_05/}: metabolic saturation figures, constraint memory time-series, interactive panel, summary JSON. All outputs are toy diagnostics.

\section*{Conflict of Interest} The author declares no conflict of interest.
\section*{Funding} This research received no external funding.
\section*{Author Contributions} Conceptualization, formal analysis, runtime engineering, and manuscript preparation: S.B.M.
\section*{Data Availability} All computational models and scripts are in the \texttt{supplementary/} directory.

\section{Reference Layer}
\bibliographystyle{plainnat}
\bibliography{afl_biology_references}

\end{document}
